If every lane privately allocates a host expert arena of size $M_E$, then with $N$ replicas the host requirement is approximately
\begin{equation}
M_{\mathrm{private}} \approx N M_E + \sum_{i=1}^{N} M_{\mathrm{lane},i},
\label{eq:private}
\end{equation}
where $M_{\mathrm{lane},i}$ covers private session, KV, runtime, and bookkeeping state. With one shared expert arena,
\begin{equation}
M_{\mathrm{shared}} \approx M_E + \sum_{i=1}^{N} M_{\mathrm{lane},i}.
\label{eq:shared}
\end{equation}
The savings approach $(N-1)M_E$ as the shared object dominates process-private state.

\subsection{State ownership}
The design shares only the large model state that is naturally identical across replicas. The shared category is the host expert backing. Lane-local state includes CUDA contexts and streams, GPU hot-expert caches, adaptive-residency metadata, GPU-resident KV, host-KV/session state, speculative/MTP state, request state, and cancellation. This boundary avoids turning normal serving progress into a distributed consistency problem. The current file-backed mapping remains writable at the OS level after initialization, so memory corruption of the shared arena is a shared failure mode; this prototype does not claim fault isolation of the weight backing.

\section{Implementation}
We implement the architecture as an opt-in compatibility layer around Strata~\cite{strata}. Each GPU lane is an ordinary Strata engine process. A supervisor sanitizes generated lane configurations so inherited \texttt{gpu} and \texttt{layer\_split} settings cannot accidentally expand one lane back into coupled multi-GPU execution.

\paragraph{Request scheduling.}
The current supervisor does not perform intra-lane continuous batching. A generation request leases one free lane, and that lane handles one active generation request at a time. If more requests arrive than there are lanes, the supervisor blocks them in a FIFO free-lane queue and dispatches each waiting request to the first lane that becomes available; the lane is returned to the pool when the request completes. This is therefore request-level queueing with first-free scheduling rather than fixed round-robin placement. The lane process itself is long-lived, so its GPU hot-expert cache and other lane-local runtime state persist across successive requests; a new request does not reconstruct the expert cache from scratch. The scaling experiments in this paper intentionally use one active request per lane and therefore do not characterize overload queueing latency, tail latency, or the throughput of $M>N$ concurrent requests.

On Linux, the fork intercepts only the exact pinned-arena allocation when the shared-arena path and byte count are explicitly supplied. Otherwise, the upstream anonymous/hugetlb allocation path remains unchanged. For the evaluated 0.1.27 baseline, the retained upstream pinned implementation is byte-identical to upstream Strata 0.1.27; the fork changes the allocation ownership boundary rather than the numerical MoE kernels. The shared mapping is \texttt{rw-s}: expert contents are logically read-mostly after population, but read-only protection is not yet enforced with \texttt{mprotect}.

The paper-validation fork commit is \texttt{6cf101d}; upstream Strata 0.1.27 is \texttt{a790805}. The evaluated binary is \texttt{build-production-027/strata}; its full commit IDs and SHA-256 are pinned in the companion artifact. The companion artifact retains repetition-level CSVs and methodology documents~\cite{artifact}.

\section{Experimental Methodology}
\begin{table}[t]
\centering
\caption{Reference host and paper-validation baseline.}
\label{tab:host}
\begin{tabular}{@{}ll@{}}
\toprule
CPU & Ryzen 9 9950X3D, 16C/32T \\
RAM & 128 GB DDR5 \\
GPU0 & RTX 5070 Ti 16 GB, PCIe x8 \\
GPU1 & RTX 5070 Ti 16 GB, PCIe x4 \\
GPU2 & RTX 5070 Ti 16 GB, PCIe x8 \\
GPU3 & RTX 5060 Ti 16 GB, PCIe x4 \\
Driver / CUDA & 615.71.09 / 13.4 \\
Engine & Strata 0.1.27 \\
